\section*{Chapitre \ref{ch:ll:pre-trade-risk-and-kill-switches} --- Contrôle de risque pré-négociation et kill switches}

\begin{solution}{exo:ll:pre-trade-risk-and-kill-switches:1}
La bande vaut $250\,000 \times 500 / 10\,000 = 12\,500$~: les achats jusqu'à 262\,500 et les ventes jusqu'à 237\,500 passent. L'achat à 26.2500
et la vente à 23.7500 passent (ils sont sur la limite)~; l'achat à 26.2600 et la vente à 23.7400 sont refusés.
\end{solution}

\begin{solution}{exo:ll:pre-trade-risk-and-kill-switches:2}
L'exposition longue dans le pire cas vaut $2\,000 + 1\,500 + 1\,000 = 4\,500$, qu'un ordre ait été acquitté ou non~; un achat de
600 porterait l'exposition à 5\,100~: refusé. Le plus grand achat accepté est de 500.
\end{solution}

\begin{solution}{exo:ll:pre-trade-risk-and-kill-switches:3}
$4 \times 10^6 / 2\,700~\text{s} \approx 1\,481$ exécutions par seconde, et $397 \times 10^6 / 4 \times 10^6 \approx 99$ actions par
exécution~: de petits ordres, et beaucoup d'entre eux.
\end{solution}

\begin{solution}{exo:ll:pre-trade-risk-and-kill-switches:4}
Un seul~: les quatorze autres reprennent le même ordre dans la fenêtre. Un emballement qui ferait varier sa quantité d'une action, ou son prix d'un tick,
d'un ordre à l'autre, passerait tous les contrôles~; le contrôle de doublons attrape les boucles qui se répètent à l'identique, pas les boucles en général.
\end{solution}

\begin{solution}{exo:ll:pre-trade-risk-and-kill-switches:5}
Jusqu'à deux secondes après le dernier \omterm{def:ll:protocols-i-fix:session}{heartbeat}, survenu au plus \qty{400}{\milli\second} avant $t$~: entre $t + 1.6$ et
$t + 2$ secondes. Une barrière qui continuerait de faire confiance à ses dernières limites jusqu'au retour du service serait à défaillance ouverte~: les limites auraient pu
avoir été resserrées, ou que le service se soit arrêté parce que quelque chose ne va pas.
\end{solution}

\begin{solution}{exo:ll:pre-trade-risk-and-kill-switches:6}
Environ 3~\% du budget. Un processus séparé sur un autre cœur demande un aller-retour par des tampons en anneau, environ \qty{200}{\nano\second}
à la médiane (chapitre~12), soit huit fois le contrôle, plus une décision sur la conduite à tenir quand il ne répond pas.
\end{solution}

\begin{solution}{exo:ll:pre-trade-risk-and-kill-switches:7}
La perte est mise à jour à chaque exécution (réalisée) et à chaque prix de référence (latente), sur le chemin des rapports de la barrière, et non dans le
contrôle~; lorsqu'elle passe sous la limite, la barrière coupe la stratégie (blocage et annulation). Ce n'est pas un contrôle sur l'ordre, puisqu'un
ordre ne modifie pas la perte~; c'est un moniteur dont l'action est une coupure. Il faut le tester avec l'emballement sur un marché en baisse.
\end{solution}

\begin{solution}{exo:ll:pre-trade-risk-and-kill-switches:8}
C'est un moniteur post-négociation, et non un contrôle de risque pré-négociation~: les ordres qu'il voit sont déjà en route, et la coupure arrive après un
saut, pendant que l'emballement continue d'envoyer (1\,500 ordres par seconde partent le temps qu'une personne lise une alerte). Si le processus s'arrête,
rien n'est contrôlé et rien ne le signale~: la conception est à défaillance ouverte. Le contrôle doit se trouver sur le chemin, avant la \omterm{def:ll:order-gateway-and-order-management:gateway}{passerelle d'ordres}~; une copie pour
la surveillance n'est qu'une seconde ligne.
\end{solution}

\begin{solution}{pb:ll:pre-trade-risk-and-kill-switches:1}
\begin{enumerate}
\item $4 \times 10^6$ exécutions en $45 \times 60 = 2\,700$ secondes~: environ 1\,481 par seconde, soit 1\,500 transactions.
\item 750\,000 actions en cinq secondes~; à 1\,500 ordres de 100 actions par seconde, $1\,500 \times 100 \times 2\,700 = 405$ millions
en 45 minutes, contre les 397 millions de l'incident.
\item 7\,500 exécutions correspondent à 749 ticks de 0.01 sur un prix de 100.00~: environ 7.5~\%.
\item Le collar est centré sur la dernière transaction, que l'emballement déplace~; chaque ordre enfant est placé cinq ticks au-dessus d'elle, bien
dans une bande de 5~\%. Le collar arrête les prix éloignés du marché, et non un marché poussé par les propres ordres de la société.
\item Chaque ordre enfant porte 100 actions, soit environ 10\,000~\$~: très en deçà de toute limite de taille ou de notionnel fixée pour les ordres ordinaires.
\item Le seau démarre plein avec 50 jetons, que les premières rafales épuisent en environ \qty{40}{\milli\second}~; ensuite,
l'emballement envoie au rythme du seau, soit 500 ordres par seconde.
\item $500 \times 100 \times 2\,700 = 135$ millions d'actions.
\item Un ordre par fenêtre~: une fois \qty{100}{\milli\second} écoulées, le même ordre n'est plus un doublon, donc environ onze
par seconde passent. Un emballement qui modifierait sa quantité d'une action à chaque fois déjouerait le contrôle.
\item 200 ordres, 20\,000 actions, avec le premier refus après \qty{130}{\milli\second}.
\item Ne compter que les exécutions ignore les \omterm{def:ll:order-gateway-and-order-management:inflight}{ordres en vol}~: la dernière rafale part avant que les exécutions de ses premiers ordres ne reviennent (un
aller-retour de \qty{70}{\micro\second} face à une rafale de \qty{15}{\micro\second}), donc dix ordres supplémentaires, 1\,000 actions, partent.
\item Avec davantage d'\omterm{def:ll:order-gateway-and-order-management:inflight}{ordres en vol}~: rafales plus grandes, ordres plus gros, cadence plus élevée, ou chemin des rapports lent (le moniteur de position de l'incident
avait pris du retard en fort volume).
\item 2.5~millions de dollars, à environ 100~\$ par action, représentent environ 25\,000 actions~; le notionnel est compté au prix de l'ordre, cinq ticks
au-dessus d'une dernière transaction en hausse, si bien que le seuil est atteint un peu plus tôt, à 24\,900.
\item \textbf{Résultat nommé.} Au rythme de l'incident, soit environ 1\,500 ordres par seconde~: sans contrôle, ou avec un collar, ou avec des limites de taille et
de notionnel par ordre, 7\,500 ordres et 750\,000 actions en cinq secondes, 405 millions d'actions transmises sur 45 minutes (l'incident~: 397 millions)~; avec une
limitation de débit à 500 par seconde, 135 millions~; avec un contrôle de doublons, 3.0~millions~; avec une limite de position
de 20\,000 comptée en vol, 200 ordres et 20\,000 actions après \qty{130}{\milli\second} (21\,000 en ne comptant que les exécutions)~; avec un
seuil de capital de 2.5~millions de dollars, 249 ordres et 24\,900 actions.
\item Le kill switch bloque la stratégie au lieu de refuser ses ordres un par un (7\,300 refus en cinq secondes sans cette mesure), annule
ses ordres ouverts et l'arrête dans chaque instrument à la fois.
\item Une limite par instrument arrête l'emballement dans chaque instrument séparément~: sur 154 instruments, 154 fois la limite. Un
seuil de capital à l'échelle de la société, ou un kill switch au premier dépassement, y est insensible.
\item La ligne du collar correspond au cap de prix à 9.5~\%~; les lignes de position sont les limites de position qui ne comptaient pas les ordres
en cours (seules les exécutions étaient comptées) et le correctif (en vol)~; la ligne du capital est le seuil à l'échelle de la société qui manquait~; le kill switch est la procédure qui manquait pour
arrêter le routeur~; les lignes de doublons et de limitation de débit sont des contrôles sur la sortie du routeur lui-même.
\item Parce que le moment où une barrière ne peut pas contrôler est celui où quelque chose ne va pas~; laisser passer les ordres n'apporte alors aucune protection au moment où
elle est nécessaire.
\item Il voit les exécutions après coup, par un chemin qui peut accuser du retard, et une personne doit agir~: à 1\,500 ordres par
seconde, chaque seconde de retard représente 150\,000 actions.
\item Des stratégies qui répètent légitimement des ordres~: un cours affiché rafraîchi au même prix et à la même taille, un ordre renvoyé après une
annulation, tranches d'un algorithme d'exécution, nouvelles tentatives d'ordres exécuter-ou-annuler.
\item Une limite sur ce que la société détient ou pourrait détenir, comptée dans le pire cas et liée à une coupure automatique.
\end{enumerate}
\end{solution}

\begin{solution}{iq:ll:pre-trade-risk-and-kill-switches:1}
Tunnel de prix (price collar), taille et notionnel maximaux par ordre, limites de position comptant les \omterm{def:ll:order-gateway-and-order-management:inflight}{ordres en vol}, seuils de capital ou de notionnel brut
au niveau du desk et de la société, cadences d'ordres et de messages, détection des doublons, nombre d'ordres ouverts, contrôles de données périmées, et un
kill switch. Ceux qui détectent un emballement sont ceux qui portent sur l'exposition accumulée (position, capital) et sur les cadences~; les contrôles par ordre
attrapent les erreurs de saisie.
\iqlookfor{la liste, et la distinction entre contrôles par ordre et contrôles d'exposition accumulée.}
\end{solution}

\begin{solution}{iq:ll:pre-trade-risk-and-kill-switches:2}
Précalculer tout (bandes du collar, limites par instrument dans des tableaux denses indexés par un petit entier), garder les expositions sous forme de
sommes courantes mises à jour par les comptes rendus, calculer chaque condition dans un masque de bits sans branchement, et refuser dès que le masque est non nul. Mesurer-la
sur les deux cas de figure.
\iqlookfor{précalcul, sommes courantes, évaluation sans branchement, mesure.}
\end{solution}

\begin{solution}{iq:ll:pre-trade-risk-and-kill-switches:3}
Elle continue de laisser passer des ordres sur les dernières limites seulement pendant leur âge maximal déclaré, puis refuse chaque ordre (fail-closed) et alerte. Les
limites ont pu être resserrées, et le silence peut être le symptôme d'une défaillance plus grave.
\iqlookfor{fail-closed, avec une période de grâce bornée.}
\end{solution}

\begin{solution}{iq:ll:pre-trade-risk-and-kill-switches:4}
Bloquer les nouveaux ordres, annuler les ordres ouverts (les annulations sont en vol jusqu'à confirmation), proposer éventuellement des ordres de soldage des positions pour
qu'une personne les confirme. La barrière elle-même, sur un dépassement franc, le desk de risque et le trader peuvent la couper~; seule une personne autorisée, après
avoir compris la cause, peut la lever.
\iqlookfor{blocage, annulation, soldage~; coupure automatique, levée manuelle.}
\end{solution}

\begin{solution}{iq:ll:pre-trade-risk-and-kill-switches:5}
Publier un instantané complet et versionné~; la barrière l'installe entre deux ordres, jamais pendant un contrôle, et enregistre la version
avec chaque décision. Les contrôles sur les nouvelles limites s'appliquent à l'exposition déjà accumulée.
\iqlookfor{instantanés atomiques entre événements, versionnés.}
\end{solution}

\begin{solution}{iq:ll:pre-trade-risk-and-kill-switches:6}
Contrôles sur la sortie de chaque composant qui génère des ordres, avant la plateforme~: limites de position et de capital dans le pire cas, liées à
une coupure automatique, contrôles de débit et de doublons, collars~; fail-closed~; surveillance post-négociation en seconde ligne. Tester en rejouant un
emballement synthétique à travers tout le chemin, un contrôle à la fois, et en vérifiant les ordres et la position auxquels chacun se déclenche.
\iqlookfor{exigences rattachées à l'incident, et rejeu d'un emballement comme test.}
\end{solution}
